% Generated draft: requires mathematical and source review.
\documentclass[11pt]{article}
\usepackage[margin=25mm]{geometry}
\usepackage{fontspec}
\setmainfont{Latin Modern Roman}
\usepackage{amsmath,amssymb,mathtools,unicode-math}
\setmathfont{Latin Modern Math}
\usepackage{enumitem,needspace,xurl,hyperref}
\setlength{\parindent}{0pt}
\setlength{\parskip}{5pt}
\setlength{\emergencystretch}{3em}
\newcommand{\N}{\mathbb{N}}
\newcommand{\Z}{\mathbb{Z}}
\newcommand{\Q}{\mathbb{Q}}
\newcommand{\R}{\mathbb{R}}
\newcommand{\C}{\mathbb{C}}
\newcommand{\F}{\mathbb{F}}
\newcommand{\A}{\mathbb{A}}
\newcommand{\PP}{\mathbb{P}}
\newcommand{\E}{\mathbb{E}}
\newcommand{\eps}{\varepsilon}
\newcommand{\dd}{\,\mathrm{d}}
\begin{document}
\section*{UnsolvedMath ID: 30001172}
\textbf{Source problem code:} catalogue entry
\subsection{Definitions and mathematical statement}
Let $j\in\N$ index the sequence of periodic traveling water waves described in the supplied record, and let $\psi^j$ denote their stream functions, with $\psi_Y^j=\partial\psi^j/\partial Y$. The coordinate $Y$, the normalization of $\psi^j$, and the fluid domains are not specified in the record. Write $D_j$ for the unspecified set over which the source takes its maximum. The stated almost-extreme condition is
\[
 m_j:=\max_{D_j}\psi_Y^j\longrightarrow 0
 \qquad (j\to\infty).
\]
This notation retains the source's maximum; it does not replace it by a supremum or by a maximum of an absolute value. The set $D_j$ and hypotheses ensuring that this maximum exists require clarification.

A stagnation point is a point where the fluid velocity relative to the traveling wave vanishes. For a planar stream-function formulation, this is equivalent to $\nabla\psi=0$. The Stokes $120^\circ$ corner condition means that two limiting free-surface branches meet with interior fluid angle $2\pi/3$, assuming such branches and their limiting tangent directions exist.

\textbf{Source-level question.} For periodic traveling water waves with general vorticity, does a sequence of almost-extreme, symmetric, monotone waves satisfying the displayed condition converge to an extreme wave satisfying the Stokes $120^\circ$ corner condition at stagnation points, rather than to a smooth wave with a bottom stagnation point?

\textbf{Unresolved formulation qualifications.} This is a question schema, not a fully specified PDE conjecture. The source extract does not define the traveling-wave equations, the admissible vorticity class, the symmetry and monotonicity conditions, the term ``extreme,'' the convergence topology, or whether passage to a subsequence is allowed. It does not state whether the assertion concerns every such sequence or a particular constructed sequence. It also does not specify which physical parameters or vorticity data remain fixed, the bed geometry, the boundary conditions, or the compactness and nondegeneracy assumptions. The corner condition is meaningful at free-surface stagnation points; the phrase ``at stagnation points'' does not identify their allowed locations or distinguish all such points from selected crests. The title refers to gravity-capillary waves, whereas the statement only says water waves and gives no surface-tension hypothesis. These omissions must be resolved before selecting a precise mathematical model or asserting the proposed conclusion.
\subsection{Short English statement}
The supplied question asks whether almost-extreme periodic waves with general vorticity approach an extreme wave with a $120^\circ$ free-surface corner, instead of a smooth wave with stagnation at the bed. The wave model, limiting procedure, and several essential hypotheses remain unspecified. The supplied status is ``partially solved''; this draft does not independently verify that assessment.
\subsection{Sources}
\begin{itemize}
\item Local project registry, permanent UnsolvedMath ID $30001172$, source problem code \texttt{catalogue entry}, title ``Extreme Gravity-Capillary Waves with General Vorticity.'' The supplied record belongs to the Oberwolfach Reports Open Problems collection and gives the proposed year as 2009; no individual proposer is supplied. Project registry locations are \texttt{data/unsolvedmath/problems.json} and the supplementary \texttt{catalog/problems.json}; the supplied extract does not identify which contains this record.
\item Supplied report citation: \emph{Wave Motion} (2009), Oberwolfach Reports, pp.~429--462, DOI 10.4171/owr/2009/07. Supplied source URL and literature-source label ``OWR Wave Motion'': \url{https://doi.org/10.4171/owr/2009/07}. The background describes Oberwolfach Reports as documenting mathematical workshops and their open-problem discussions.
\item Supplied literature-source label ``Varvaruca--Weiss Stokes conjecture with vorticity'': \url{https://doi.org/10.1016/J.ANIHPC.2012.05.001}.
\item Supplied literature-source label ``extreme waves and Stokes conjecture with vorticity'': \url{https://arxiv.org/abs/0707.2224}.
\item Supplied literature assessment, checked at the exact date \texttt{2026-08-21}: status ``partially solved.'' The curation states that the $120$-degree Stokes-corner alternative and substantial local classification for waves with vorticity were proved, while existence and exclusion of all alternative singular profiles for arbitrary vorticity are not completely settled. This is supplied dated curation derived from the OWR material, not an independently verified literature assessment; its time-sensitive status and cited claims require independent verification.
\item The registry labels the statement \texttt{corrected\_verified} and reports that its mathematical content was checked against the official Oberwolfach report, with notation, context, and question form normalized without changing the source-level claim. This verification claim is supplied metadata; the present draft has not independently checked the report or the linked literature.
\end{itemize}
\end{document}
